% Back cover of cascade_booklet: Practice Framework (portrait letter).
\documentclass[letterpaper]{article}
\usepackage{fontspec}
\usepackage{geometry}
\usepackage{xcolor}
\usepackage{colortbl}
\usepackage{array}
\usepackage{enumitem}

\geometry{letterpaper, margin=10mm}
\setmainfont{DejaVu Sans}

\definecolor{headerbg}{HTML}{2C3E50}
\definecolor{white}{HTML}{FFFFFF}
\definecolor{loopClr}{HTML}{2E7D32}
\definecolor{ratioClr}{HTML}{3F51B5}
\definecolor{axesClr}{HTML}{F57C00}
\definecolor{orderClr}{HTML}{C62828}

\newcommand{\sect}[3]{%
  \noindent\begin{minipage}{\linewidth}%
  \colorbox{#1}{\color{white}\textbf{\ #2\ }}\par\vspace{1mm}
  #3
  \end{minipage}\par\vspace{3mm}%
}

\setlist[itemize]{leftmargin=5mm, itemsep=1.5pt, topsep=2pt, parsep=0pt, partopsep=0pt}
\setlength{\parindent}{0pt}
\setlength{\parskip}{0pt}
\pagestyle{empty}

\begin{document}

\begin{center}
{\fontsize{18}{22}\selectfont\textbf{Practice Framework}}\\[-1mm]
{\small\itshape How to actually improve, every session}
\end{center}
\vspace{2mm}
\hrule
\vspace{3mm}

\fontsize{10}{13}\selectfont

\sect{loopClr}{The Feedback Loop --- every single rep}{%
No rep is wasted if you treat each one as an experiment.
\begin{itemize}
  \item \textbf{Plan}: before playing, set a specific target. Not ``play it better'' --- something like ``keep the LH thumb on the string until the RH note starts.''
  \item \textbf{Play}: execute the target honestly, at whatever tempo lets you think.
  \item \textbf{Evaluate}: did you hit the target? Be strict. 90\% is success; 50\% is information.
  \item \textbf{Adjust}: change one variable before the next rep --- slow down, shorten the phrase, or pick a narrower target. Never repeat blindly.
\end{itemize}
Liszt: ``\textit{When practicing, think ten times and play once.}''
}

\sect{ratioClr}{The 90/10 Principle}{%
To nail a skill at 100\%, stop grinding it head-on. Push a \textbf{harder related skill} to 90\% instead --- the easier one snaps into place automatically.

\vspace{1mm}
\textit{Examples for harp:}
\begin{itemize}
  \item Want clean arpeggios at tempo? Practice them with finger-numbers said out loud at 70\% tempo. When you stop speaking, tempo goes up.
  \item Want a smooth LH/RH hand-off? Do the cascade in reverse (descending then ascending). When you go forward again it feels easy.
  \item Want even dynamics? Practice crescendos and decrescendos over the same notes. Then ``flat'' feels natural.
\end{itemize}
You're not doing the target skill --- you're doing \emph{something harder}. When you remove the extra layer, the target skill runs on less effort.
}

\sect{axesClr}{Five Axes of Difficulty}{%
When ``harder'' is vague, pick one of five axes. Raise one, lower another, cycle.

\vspace{1mm}
\renewcommand{\arraystretch}{1.2}
\setlength{\tabcolsep}{4pt}
\begin{tabular}{@{}p{22mm}p{\dimexpr\linewidth-28mm\relax}@{}}
\textbf{Length} & one measure $\to$ two $\to$ phrase $\to$ whole piece. Endurance \textit{is} a skill. \\
\textbf{Complexity} & one hand $\to$ two hands; melody alone $\to$ melody + chords; add a cascade across a bar line. \\
\textbf{Depth} & not just notes: phrasing, dynamics, articulation, tone. The ``musicality'' layer. \\
\textbf{Speed} & only after the first four feel solid. Speed is the easiest axis and the most overused. \\
\textbf{Mental endurance} & how long you stay focused \textit{inside} the playing. Sustained attention, not just physical stamina. \\
\end{tabular}
}

\sect{orderClr}{3-4-2-1 Session Order}{%
Your brain's focus falls during any session. Front-load the hardest work.

\vspace{1mm}
Rank every chunk of material into four levels:
\begin{itemize}
  \item \textbf{Level 4} --- ``impossible right now.'' Can't play it cleanly even once.
  \item \textbf{Level 3} --- ``hard.'' You have to be at your A-game.
  \item \textbf{Level 2} --- ``almost there.'' Mostly clean, occasional slip.
  \item \textbf{Level 1} --- ``automatic.'' Plays itself.
\end{itemize}

\textbf{Session order: 3 $\to$ 4 $\to$ 2 $\to$ 1.}
Start at level 3 as a warm-up that demands real focus (skip level-1 warm-ups --- they set your brain to cruise). Move to level 4 while you're sharpest. Drop to level 2 as attention fades. End with level 1 for review, while your brain is already in ``just play'' mode.

\textit{This works at any skill level --- your level 4 as a beginner and a professional's level 4 are different pieces, but the structure is identical.}
}

\vspace{2mm}
\hrule
\vspace{1mm}
\small\centering\itshape
Chopin told his students two hours a day was enough. They learned his etudes. The difference was never quantity.

\end{document}
